\songtitle{The Secret Fires}

Music theater version.
\tag*{2023}\tag{original-seed}\tag{social-media-seed}\tag{musical-theater}

\subsection{Notes}
Check notes to \emph{The Secret Fires} Ballad Pop version.\\
This is based on a second iteration by Gemini, asking for a faster pace, and fed that 30s product to Suno.

Lyrics outro tweaked manually .

\subsection{Style}
Musical theater showtune with a swing-influenced jazz pop arrangement. %
A baritone male vocal performs with clear diction and theatrical phrasing. %
The instrumentation features a bright acoustic piano playing syncopated block chords and melodic fills, an upright bass playing a walking bassline, and a drum kit using brushes on the snare with a steady swing rhythm. %
A brass section consisting of trumpets and trombones provides staccato stabs and melodic counterpoints. %
The tempo is 128 BPM in 4/4 time, likely in the key of C major.

\begin{lyrics}
\annotation{baritone male vocals}\\
I take my eggs without a salt,\\
My broth without the green,\\
My coffee dark, my fish un-tart,\\
The plainest you have seen.

\songpart{Verse 2}
They call me bland, they mock my taste,\\
And think my world is gray,\\
\annotation{brass section enters with staccato stabs}\\
But they miss the secret fires I light\\
When they have looked away.

\songpart{Verse 3}
For they don't see the heavy spice,\\
The oils that richly gleam,\\
The pepper on the chicken rice,\\
The vinegar in steam.

\songpart{Verse 4}
A hidden warmth within the dish,\\
A sudden, sharp delight—\\
They think my life is color-blind,\\
But I feast in the night.

\songpart{Bridge}
the textures and the flavors\\
the bitterness\\
delicacies of the mix

\songpart{Outro}
They call me bland, they mock my taste,\\
And think my world is gray,\\
They think my life is color-blind,\\
But I feast in the night.

But I feast in the night.
\end{lyrics}

